% Generated draft: requires mathematical and source review.
\documentclass[11pt]{article}
\usepackage[margin=25mm]{geometry}
\usepackage{fontspec}
\setmainfont{Latin Modern Roman}
\usepackage{amsmath,amssymb,mathtools,unicode-math}
\setmathfont{Latin Modern Math}
\usepackage{enumitem,needspace,xurl,hyperref}
\setlength{\parindent}{0pt}
\setlength{\parskip}{5pt}
\setlength{\emergencystretch}{3em}
\newcommand{\N}{\mathbb{N}}
\newcommand{\Z}{\mathbb{Z}}
\newcommand{\Q}{\mathbb{Q}}
\newcommand{\R}{\mathbb{R}}
\newcommand{\C}{\mathbb{C}}
\newcommand{\F}{\mathbb{F}}
\newcommand{\A}{\mathbb{A}}
\newcommand{\PP}{\mathbb{P}}
\newcommand{\E}{\mathbb{E}}
\newcommand{\eps}{\varepsilon}
\newcommand{\dd}{\,\mathrm{d}}
\begin{document}
\section*{UnsolvedMath ID: 30002506}
\textbf{Source problem code:} catalogue entry
\subsection{Definitions and mathematical statement}
The supplied question concerns a Banach sequence space \(X\). Write \(X_n\) for its first \(n\) coordinate space and \(B_{X_n}\) for its open unit ball. The exact convention for a Banach sequence space and the construction of \(X_n\) are not specified in the supplied record. Under the inherited coordinate-subspace interpretation, if the coordinate vectors \(e_j\) belong to \(X\), then
\[
X_n=\operatorname{span}\{e_1,\ldots,e_n\},\qquad
\|(z_1,\ldots,z_n)\|_{X_n}
=\|(z_1,\ldots,z_n,0,\ldots)\|_X,
\]
with \(B_{X_n}=\{z\in X_n:\|z\|_{X_n}<1\}\).

For clarity, a usual complex holomorphic interpretation of the Bohr radius is the following. Let \(H^\infty(B_{X_n})\) be the bounded holomorphic functions on \(B_{X_n}\), with
\[
\|f\|_\infty=\sup_{z\in B_{X_n}}|f(z)|.
\]
For \(\alpha\in\mathbb Z_{\geq0}^n\), put \(z^\alpha=\prod_{j=1}^n z_j^{\alpha_j}\), and let \(a_\alpha(f)\) be the Taylor coefficients of \(f\) at the origin. Under this convention,
\[
K(B_{X_n})=
\sup\left\{r\in[0,1]:
\begin{array}{l}
\displaystyle\sum_{\alpha\in\mathbb Z_{\geq0}^n}
|a_\alpha(f)z^\alpha|\leq\|f\|_\infty\\[2pt]
\text{for every }f\in H^\infty(B_{X_n})
\text{ and every }z\in rB_{X_n}
\end{array}\right\},
\]
where \(rB_{X_n}=\{rw:w\in B_{X_n}\}\), including \(0B_{X_n}=\{0\}\). This definition makes the customary interpretation explicit; the supplied record does not itself specify the scalar field or the defining Bohr inequality, so its agreement with the source convention requires review.

In the complex interpretation, define
\[
\ell^1=\left\{x=(x_j)_{j\geq1}\in\C^{\{1,2,\ldots\}}:
\sum_{j=1}^{\infty}|x_j|<\infty\right\}.
\]
The supplied mathematical question is whether, for every Banach sequence space \(X\) in the intended source class,
\[
\inf_{n\geq1}K(B_{X_n})>0
\quad\Longrightarrow\quad
X=\ell^1\ \text{as sets}.
\]
The hypothesis means that there exists \(\delta>0\) such that \(K(B_{X_n})\geq\delta\) for every integer \(n\geq1\). The conclusion asks for equality of the underlying sequence sets, not equality of norms. The intended scalar field, sequence-space axioms, coordinate-space construction, and Bohr-radius convention remain source qualifications to be checked; no additional symmetry, solidity, or density hypothesis is asserted here.
\subsection{Short English statement}
Must every Banach sequence space whose finite coordinate unit balls have Bohr radii bounded below by one positive constant consist exactly of the absolutely summable sequences? Equality is requested only as sets. The supplied record labels this question open; the precise sequence-space and Bohr-radius conventions require review.
\subsection{Sources}
\begin{itemize}
\item Supplied local UnsolvedMath registry record, permanent ID \(30002506\), source problem code \texttt{catalogue entry}, title ``Bohr-Radius Characterization of \(\ell^1\)'', in the Oberwolfach Reports Open Problems collection. The project registry locations are \texttt{data/unsolvedmath/problems.json} and \texttt{catalog/problems.json}; the supplied record is the evidence used here.
\item Supplied report citation: \emph{Dirichlet Series and Function Theory in Polydiscs} (2014), Oberwolfach Reports, DOI \texttt{10.4171/owr/2014/06}, pp.\ 335--393. \url{https://doi.org/10.4171/owr/2014/06}. The record attributes the problem to the report's workshop open-problem discussions, gives proposed year 2014, and identifies no individual proposer or exact problem page.
\item Supplied literature source labelled ``Recent monomial-convergence work'': \url{https://riunet.upv.es/bitstreams/aaf75982-880b-4910-a99e-76252cefcfad/download}. No author, title, or publication details for this link are supplied.
\item Literature assessment dated exactly \texttt{2026-08-21}: the supplied curation labels the problem open and reports that no general proof or counterexample was found for the asserted implication, while available work treats related Bohr-radius and monomial-convergence questions. The listed checked sources are the report and the monomial-convergence link above. This is supplied dated literature triage derived from the OWR curation, not an independently verified assessment; time-sensitive status and cited claims require independent verification.
\item The supplied statement status is \texttt{corrected\_verified}, with a supplied verification note claiming comparison against the official report and preservation of the hypotheses, notation, and open target. That verification claim is retained as registry provenance and was not independently checked for this draft.
\end{itemize}
\end{document}
